% ===========================================================================
\section{Faces of the transform}\label{sec:faces}
\warmth{0}\quad\whisper{Fast computation, sampling, filtering, PDEs: one dictionary, cashed out.}

\begin{strand}
Because Fourier diagonalizes every shift-invariant operation, a long list of classical problems
collapse to ``transform, multiply, transform back''. Here is the keyring.
\end{strand}

\block{The keyring}
\noindent\begin{tabularx}{\linewidth}{@{}l L@{}}
\toprule
{\color{indigo}face} & \multicolumn{1}{l}{\color{indigo}the idea}\\
\midrule
the \keyword{FFT}              & compute the length-$N$ DFT in $O(N\log N)$, not $O(N^2)$, by splitting even/odd indices and recursing\\
\keyword{sampling} (Nyquist)   & a signal band-limited to $\abs\xi\le B$ is fully determined by samples at rate $\ge 2B$\\
\keyword{filtering}            & multiply $\widehat f$ by a mask $=$ convolve $f$ with the kernel; low-pass, high-pass, denoise\\
solving \keyword{PDEs}         & the heat equation $u_t=u_{xx}$ becomes $\widehat u_t=-(2\pi\xi)^2\widehat u$: an ODE \emph{per frequency}\\
\keyword{circulant} matrices   & all diagonalized by the Fourier matrix; eigenvalues are the $\widehat{}$'s\\
\bottomrule
\end{tabularx}

\begin{theorem}[the FFT, in one line]\label{thm:fft}
Splitting a length-$N$ DFT into its even- and odd-indexed halves gives
$\widehat x_k=\widehat{(\text{even})}_k+e^{-2\pi i k/N}\,\widehat{(\text{odd})}_k$, two DFTs of size
$N/2$ plus $O(N)$ glue. Recursing costs $T(N)=2T(N/2)+O(N)=O(N\log N)$.
\end{theorem}

\begin{yourturn}
The heat equation, diagonalized. Transforming $u_t=u_{xx}$ in $x$ turns it into a simple ODE in
$t$ for each frequency:
\[
  \widehat u_t(\xi,t)=-(2\pi\xi)^2\,\widehat u(\xi,t)\quad\Longrightarrow\quad
  \widehat u(\xi,t)=\fillin[3.5cm].
\]
High frequencies therefore decay \fillin[2cm] than low ones. (That is exactly why heat smooths.)
\recall{$\widehat u(\xi,t)=\widehat u(\xi,0)\,e^{-(2\pi\xi)^2 t}$; high frequencies decay \emph{faster} (the $\xi^2$ in the exponent). Smoothing $=$ killing high frequencies.}
\end{yourturn}

\block{The one idea, recovered}
\begin{strand}
Every face above is the same sentence: \keyword{Fourier diagonalizes shift}. Convolution,
differentiation, circulant matrices, and constant-coefficient PDEs are all built from translation,
so all of them become multiplication once you change to the frequency basis. Read the master table
again; you can now fill, and explain, every row.
\end{strand}

\loose{Where to go next: replace $\R$ by a finite group and you get the DFT and group representations; demand localization in \emph{both} domains and you get wavelets; take the group to be $\{0,1\}^n$ and you get the \keyword{Fourier analysis of Boolean functions} that powers much of theoretical CS.}
